% ===== §8 The formal-modelling dilemma =====
\section{The Formal Modelling of the Value Cycle, and the Question of the Other}
\label{sec:modelling}

\noindent
The quantum-structural grammar of \xref{sec:grammar} articulates participatory generation, but it leaves unaddressed a difficulty that conditions the entire formal programme. The difficulty is that the subject who creates value is reconfigured, within the cycle of value, by the value it creates (\xref{sec:reflexive}); the generating mechanism is rewritten, in the process of generation, by what the process generates. This section establishes the precise character of the difficulty. It shows that the difficulty is not, as it first appears, the inadequacy of the standard formal tools, which can in fact be made adequate; that every device by which they are made adequate shares a single presupposition; that this presupposition is the existence of a fixed underlying structure, a metalanguage or big Other of the formal system; and that the difficulty is therefore identical to the question, posed throughout the series, of whether there is an Other of the Other. The section closes by identifying the one class of formal tools structurally compatible with a negative answer.

\subsection{The unified form of the difficulty}
\label{sec:mod-form}

The difficulty assumes one form across three modelling vocabularies. In the vocabulary of generative grammar, the set of production rules is not fixed but is altered, during generation, by what the grammar generates: a self-modifying grammar. In the vocabulary of dynamical systems, the law of evolution is not fixed but is altered, during evolution, by the states the system produces: where a standard system evolves as $\dot{x} = f(x)$ with $f$ fixed, the present case requires that $f$ itself evolve under the influence of the phenomena its evolution produces. In the vocabulary of field theory, the coupling structure among the constituents is not fixed but is reconfigured, during the dynamics, by the emergent the coupling produces (\xref{sec:registers}). The three are one difficulty: the reconfiguration of the generating structure by the generated phenomenon, the reflexive structure of \xref{sec:reflexive} posed as a problem of formal modelling.

\subsection{The standard tools can be made adequate}
\label{sec:mod-rescue}

It is tempting to conclude that the standard formal tools fail before this difficulty. They do not, and the supposition that they do must be set aside before the genuine difficulty can be located. For in each vocabulary the variation of the generating structure can be accommodated by the introduction of a fixed higher-order structure that produces the variation while itself remaining invariant.

In the grammatical vocabulary, a self-modifying grammar whose rule-set $G$ is altered during generation can be modelled by a fixed meta-grammar $M$ that produces and rewrites $G$: $M$ does not vary, and the variation of $G$ is the output of the invariant $M$. In the dynamical vocabulary, a system whose law $f$ varies can be modelled by a fixed higher-order law: $\dot{x} = f(x, p)$, $\dot{p} = g(x, p)$, where the variation of $f$ is governed by the invariant pair $(f, g)$, or, if $g$ too must vary, by a fixed $(f, g, h, \dots)$ at some order. In each case the resulting system remains within the recursively enumerable: a fixed Turing machine carrying the meta-rules as part of its invariant program simulates the self-modification, and no super-recursive capacity is required or claimed.

\begin{claim}[The standard tools are rescued by a fixed higher-order structure]
The variation of the generating structure does not place the value cycle beyond the standard formal tools. In each vocabulary the variation is accommodated by a fixed higher-order structure, a meta-grammar or a higher-order law, that produces the variation while remaining itself invariant, and the resulting system remains within the recursively enumerable. The standard tools do not fail; they are made adequate by the introduction of an invariant higher-order structure. The supposition that the value cycle exceeds the formal power of the standard tools is therefore to be rejected. The genuine difficulty lies elsewhere: not in the power of the tools, but in what their rescue presupposes.
\end{claim}

\subsection{The shared presupposition, and the identity of the field-theoretic and measurement rescues}
\label{sec:mod-shared}

Every rescue of the preceding kind shares one presupposition: a fixed underlying structure, invariant beneath the variation it produces. The meta-grammar $M$ is fixed; the higher-order law $(f, g, \dots)$ is fixed; and the rescue consists precisely in locating an invariant structure of which the observed variation is the determinate product. This holds equally of the two rescues that appear, at first, to differ in kind. The field-theoretic reformulation, in which the coupling among the registers is itself a dynamical variable, alters nothing: it translates the difficulty into the vocabulary of variable coupling without removing it, and its own rescue is again a fixed higher-order law governing the variation of the coupling. And the measurement rescue, in which the underlying field-theoretic dynamics is held fixed and the entire appearance of reconfiguration is located in the structure of measurement (the participatory measurement of \xref{sec:measurement}), is a special case of the fixed-meta-structure rescue: it holds a fixed underlying dynamics, the analogue of the fixed $M$, and locates the variation in the derived relation of observation, the analogue of $M$'s production of $G$. The measurement device introduces no formal capacity the meta-grammar lacks; it relocates the variation to the observational relation, but it relocates it onto a fixed underlying structure exactly as the meta-grammar does.

\begin{claim}[The measurement rescue is a special case of the fixed-meta-structure rescue]
Holding the underlying dynamics fixed and locating reconfiguration in measurement does not differ, in its formal presupposition, from holding a meta-grammar fixed and locating variation in its production of the rule-set. Both presuppose a fixed underlying structure of which the observed variation is the derived product; the measurement rescue is the measurement-theoretic instance of the general fixed-meta-structure rescue. The introduction of an observation $O$ that registers the emergence of value and the evolution of the subject solves no formal problem that the fixed meta-grammar does not already solve, and incurs the same presupposition: a fixed underlying structure beneath the reflexive appearances.
\end{claim}

\subsection{The presupposed structure as metalanguage, and the regress of the Other}
\label{sec:mod-metalanguage}

The fixed underlying structure that every rescue presupposes is to be named. It is the invariant generative law of which every actual value cycle, and every reconfiguration of a subject within it, is a determinate product. Call it the \emph{meta-value-grammar}: the fixed generative structure beneath the reflexive appearances, the meta-grammar $M$ in the grammatical vocabulary, the higher-order law in the dynamical, the fixed underlying field in the measurement rescue. Every formal rescue of the value cycle's modelability presupposes the meta-value-grammar; the question of whether the value cycle is formally modelable, by a fixed structure, is identical to the question of whether the meta-value-grammar exists.

This identifies the difficulty with the question the series has pursued under another name. The meta-value-grammar is a metalanguage in the precise sense of \xref{sec:breaking}: a fixed structure standing beneath, and accounting for, the phenomena, from which their variation is surveyed and produced. And it is the big Other in the precise sense of the series: the fixed locus that would guarantee the determinacy of the value cycle from outside its reflexive movement. The difficulty of formal modelling is therefore not a technical difficulty separable from the philosophical commitments of the paper; it is those commitments, posed in the formal register. To ask whether the value cycle admits a fixed formal model is to ask whether there is a metalanguage, whether there is a big Other.

And the regress is the one the series has identified. The meta-grammar $M$ rescues the modelability of the varying rule-set $G$, but $M$ in the present case is no more fixed than $G$: the generative structure of intimacy is itself reconfigured across the history of a relation, so that $M$ requires its own meta-grammar $M'$, and $M'$ its own, in a regress that does not terminate. This is, in the formal register, the regress Lacan identifies and arrests with the thesis that there is no Other of the Other: the big Other is itself lacking, barred, without a further Other to guarantee it \citep{lacan_no_other,fink1995,evans1996}. The formal correlate is exact: there is no terminal fixed meta-value-grammar, since each presupposed fixed structure requires a further fixed structure to produce its own variation, and the regress of meta-grammars does not close. The rescue of modelability by a fixed higher-order structure does not eliminate the difficulty; it displaces it upward by one level and reproduces it there.

\begin{claim}[The modelling dilemma is the question of the Other of the Other]
The question whether the reflexive value cycle admits a fixed formal model is identical to the question whether there exists a fixed meta-value-grammar, which is the question whether there is a metalanguage and a big Other. Every fixed higher-order rescue presupposes such a structure and thereby incurs the regress of the Other: the presupposed fixed structure is itself variable and requires a further fixed structure, without termination. The formal dilemma is therefore not separable from the central commitment of the series. To affirm a fixed formal model of the value cycle is to affirm an Other of the Other; to deny the Other of the Other, with Lacan and with the ontology of \xref{sec:onto-refuses}, is to deny that the value cycle admits a fixed formal model of the standard kind.
\end{claim}

\subsection{Background-independent generativity as the form of the barred Other}
\label{sec:mod-gft}

The dilemma so stated appears to permit only two responses: the renunciation of formal modelling, or the affirmation of a fixed meta-value-grammar that the series elsewhere denies. But the formal situation admits a third possibility, and it is the one structurally compatible with the denial of the Other of the Other. The standard formal tools, including ordinary quantum field theory on a fixed background, presuppose a fixed underlying structure: ordinary quantum field theory is defined on a fixed spacetime, the invariant stage beneath the field's dynamics, and is in this respect the most explicit form of the fixed meta-value-grammar. But a class of formal frameworks has been developed precisely to remove the fixed background: the background-independent approaches to quantum gravity, and among them group field theory, in which the field is defined not on a fixed spacetime but on a group manifold, and spacetime and geometry are not presupposed but emerge as collective phenomena of the field's interaction \citep{oriti2012,oriti2014,rovelli2004}.

The pertinence of this class of frameworks is structural and is offered as such, not as a claim that the value cycle is a group field theory. What background-independent generativity supplies is a formal framework in which the underlying structure is itself emergent rather than fixed: there is no invariant background stage beneath the dynamics, the stage itself being a product of the generativity rather than its presupposition. This is the formal correlate of the barred Other. The denial of the Other of the Other is not the denial that there is any generativity, but the thesis that the generativity is itself lacking, without a fixed further structure to ground it; and a background-independent framework realises exactly this, a generativity whose own ground is emergent and not presupposed.

\begin{claim}[Background-independent generativity as the form compatible with the barred Other]
A background-independent formal framework, in which the underlying structure is itself emergent rather than a fixed presupposed background, is the form of formal tool structurally compatible with the denial of the Other of the Other. The meta-value-grammar it would supply is not a fixed underlying structure, the unbarred Other, but a generativity whose own ground is emergent and lacking, the barred Other. Group field theory and the background-independent approaches to quantum gravity are advanced here as the structural exemplars of this possibility, and not as a physical model of intimacy: they exhibit, in a developed formalism, a generativity that does not presuppose a fixed background, and thereby indicate the one direction in which the formal modelling of the reflexive value cycle might proceed without positing the fixed meta-value-grammar the series denies.
\end{claim}

\subsection{The dilemma restated}
\label{sec:mod-dilemma}

The dilemma is now stated in its proper form. It is not the technical dilemma of whether the standard tools are powerful enough, for they are; it is the philosophical dilemma, posed in the formal register, of whether there is an Other of the Other. Three positions remain, and the paper does not adjudicate among them, though the third is now a determinate formal programme rather than a suspension.

The first position affirms a fixed meta-value-grammar: it accepts a fixed underlying structure, of which the reflexive appearances are determinate products, and thereby accepts an Other of the Other, at the cost of the central commitment of the series. It is the position of the standard formal rescue, and its cost is the abandonment of the denial of the metalanguage.

The second position renounces formal modelling: it holds that the generativity is irreducibly without a fixed ground, that no fixed meta-value-grammar exists, and that the value cycle is therefore not modelable by any fixed structure and is to be witnessed rather than modelled (\xref{sec:walk}). It is the most direct fidelity to the denial of the Other of the Other, at the cost of conceding that the central object of the paper exceeds fixed formal representation.

The third position seeks a background-independent formalism: it holds that the denial of the Other of the Other does not require the renunciation of formalism as such, but the renunciation of the fixed background, and that a formal framework whose own underlying structure is emergent, on the model of group field theory, is the form a non-settling formalisation of the value cycle would take. This position does not resolve the dilemma, since whether such a formalism can be constructed for the value cycle is not here established; but it converts the third position from a suspension into a programme, and identifies the formal direction, background-independent generativity as the form of the barred Other, in which the matter is to be pursued. This is the point at which the formal programme of the paper hands forward to the Value-Foam programme and to future work, and the point at which the formal dilemma is seen to be, in its entirety, the central question of the series posed in the register of formal modelling.
